% Last Updated: March 1, 2024
\documentclass[0-main]{subfiles}

\begin{document}
\ifSubfilesClassLoaded{
	% \title{\tbf{Untitled}}
	% \author{Cong Wen, xx/20xx}
	% \maketitle
	\tableofcontents
}{}
% ----------------------------------------------------------------------------------------------------


\ssc{On number fields}

Memorandum of number fields.

\sss{Quadratic fields}

\begin{thm}
		Quadratic fields are of the form $K=\bQ(\sqrt{D})$, where $D\neq1$ is a square-free integer. Then
		\[
				\Dta_{K}=\begin{cases}
								D & D\eqv1\bmod{4} \\
								4D & D\eqv2,3\bmod{4}
							\end{cases},\quad \cO_{K}=\bZ\sS{\frac{1}{2}(\Dta_{K}+\sqrt{\Dta_{K}})}.
		\]
		The splitting type of any prime $p$ in $K$ is given by
		\[
				\text{$p$ ramified}\aolR p\mid \Dta_{K},\quad \text{$p$ split}\aolR \begin{cases}
																																		\Dta_{K}\eqv1\bmod{8} & p=2 \\
																																		(\Dta_{K}/p)=1 & p\neq2
																 \end{cases},\quad \text{$p$ inert}\aolR \begin{cases}
																									   \Dta_{K}\eqv 5\bmod{8} & p=2 \\
																									   (\Dta_{K}/p)=-1 & p\neq2
																									 \end{cases}.
		\]
\end{thm}

% \begin{thm}
% 		Let $N\gqs1$ be an integer, then for any prime $p\nmid 2N$,
% 		\[
% 				\xis x,y\in\bZ, p=x^{2}+Ny^{2} \aolR \text{$(-N/p)=1$ and $H_{-4N}(x)\eqv0\bmod{p}$ has an integer solution}.
% 		\]
% \end{thm}


\begin{qst}
	Let $\cO\sbs K$ be an imaginary quadratic order of discriminant $D$, $H=K(j_{D})$ the ring class field, $F$ the genus field, denote by $\tau$ the complex conjugation, then one has the field tower,
	\[
		\begin{tikzcd}
			K\ar[r, dash]\ar[d, dash, "\sA{\tau}"] & F\ar[r, dash, "\Cl(\cO)^{2}"]\ar[d, dash, "\sA{\tau}"] & H\ar[d, dash, "\sA{\tau}"] \\
			\bQ\ar[r, dash] & F\xP\ar[r, dash] & H\xP
		\end{tikzcd}
	\]
	\begin{itm}
		\item[\TODO] Determine the splitting type of any prime $p$ in the field tower.
		\item[\TODO] Determine the factorization of $H_{D}(x)$ modulo any prime $p$. (Use also 2108.00168)
	\end{itm}
\end{qst}


% \blue{ - - - progress bar - - - }

Do continued fractions, fundamental units etc...

% See also \cite{Cohen2000}, also check out \cite[Section~6.2]{Cohen1993}



\sss{Biquadratic fields}
Ref:~\cite{Williams1970}.

% % \begin{thm}
% % 		Biquadratic fields are of the form $K=\bQ(\sqrt{ml},\sqrt{nl},\sqrt{mn})$, where $m,n,l\neq1$ are coprime square-free integers. Then
% % \end{thm}

% Biquadratic fields are of the form $K=\bQ(\sqrt{m},\sqrt{n})$, where $m,n$ are distinct square-free integers not equal to $1$, write $m=lm_{1},n=ln_{1}$ such that $(m_{1},n_{1})=1$. Then
% \begin{align*}
%   \cO_{K}&=\begin{cases}
% 			  \bZ+\bZ\frac{1+\sqrt{m}}{2}+\bZ\frac{1+\sqrt{n}}{2}+\bZ\frac{1+\sqrt{m}+\sqrt{n}+\sqrt{m_{1}n_{1}}}{2} & (m,n)\eqv(1,1)\bmod{4},\quad (m_{1},n_{1})\eqv(1,1)\bmod{4} \\
% 			  \bZ+\bZ\frac{1+\sqrt{m}}{2}+\bZ\frac{1+\sqrt{n}}{2}+\bZ\frac{1-\sqrt{m}+\sqrt{n}+\sqrt{m_{1}n_{1}}}{2} & (m,n)\eqv(1,1)\bmod{4},\quad (m_{1},n_{1})\eqv(3,3)\bmod{4} \\
% 			  \bZ+\bZ\frac{1+\sqrt{m}}{2}+\bZ\sqrt{n}+\bZ\frac{\sqrt{n}+\sqrt{m_{1}n_{1}}}{2} & (m,n)\eqv(1,2)\bmod{4} \\
% 			  \bZ+\bZ\sqrt{m}+\bZ\sqrt{n}+\bZ\frac{\sqrt{m}+\sqrt{m_{1}n_{1}}}{2} & (m,n)\eqv(2,3)\bmod{4} \\
% 			  \bZ+\bZ\sqrt{m}+\bZ\frac{\sqrt{m}+\sqrt{n}}{2}+\bZ\frac{1+\sqrt{m_{1}n_{1}}}{2} & (m,n)\eqv(3,3)\bmod{4}
% 			\end{cases} \\
%   \Dta_{K}&=\begin{cases}
% 		  m^{2}n^{2} & (m,n)\eqv(1,1)\bmod{4} \\
% 		  16m^{2}n^{2} & (m,n)\eqv(1,2), (3,3)\bmod{4} \\
% 		  64m^{2}n^{2} & (m,n)\eqv(2,3) \bmod{4}
% 		\end{cases}.
% \end{align*}

% \red{(1)TODO}

% The splitting type of any rational prime $p$ in $K$ is given by

\sss{Cubic fields}
Ref:~\cite{MW2019}


\sss{Quartic fields}
Ref:~\cite{HY2012}


\sss{Cyclotomic fields}
Ref:~\cite{Washington1997}.

% Cyclotomic fields are of the form $K_{N}=\bQ(\zta_{N})$, where $N\gqs3$ is an integer. Now assume that $N=p^{r}$ below.

% The discriminant, ring of integers are given by
% \[
% 	\Dta_{K}=\frac{(-1)^{\phi(N)/2}N^{\phi(N)}}{\prod_{p\mid N}p^{\phi(N)/(p-1)}},\quad \cO_{K}=\bZ[\zta_{N}].
% \]

% The splitting type of any rational prime $p$ in $K_{N}$ is given by

% If $p=2, r\gqs3$, then

% If $p\gqs3, r\gqs2$, then

% % \[
% % 	\text{$p$ ramified}\aolR p\eqv0\bmod{N},\quad \text{$p$ split}\aolR p\eqv1\bmod{N},\quad \text{$p$ inert}\aolR \text{$p$ is primitive modulo $N$}.
% % \]


% $K_{N}$ has a maximal real subfield $K_{N}\xP=\bQ(\zta_{N}+\zta_{N}\xI)$
% \[
% 	\Dta_{K_{N}\xP}=???,\quad \cO_{K_{N}\xP}=\bZ[\zta_{N}+\zta_{N}\xI].
% \]
% The splitting type of any rational prime $p$ in $K_{N}\xP$ is given by

% \red{(3)TODO}
% Can divide into three parts, $K_{2^{r}}\xP$ for $r\gqs3$, $K_{p^{r}}$ for $p\gqs3, r\gqs2$, and in general

% $K_{p^{r}}/K_{p^{r}}\xP$ is ramified only at primes above $p$ and $\ift$.

% If $N$ is not a prime power, then $K_{N}/K_{N}\xP$ is ramified only at $\ift$.

% $K_{p},p>2$ has a unique quadratic subfield $\bQ(\sqrt{(-1)^{(p-1)/2}p})$.




% \sct{On quaternion algebras}
% Can one classify all maximal orders in $B_{p,\ift}$?

% Use \cite{Voight2021}.



% \sct{On finite groups}

% \ssc{Maximal subgroups of $\GL_{2}(\bF_{p})$}
% Use \cite[Section~2]{Serre1972}, describe and classify all its subgroups.

% ----------------------------------------------------------------------------------------------------
\ifSubfilesClassLoaded{
	\bibliographystyle{amsalpha}
	\bibliography{/Users/wenc/Documents/Math/universal-ref}
}{}
\end{document}

